\documentclass[aspectratio=169]{beamer}

% =====================================================================
% ART DECK 05 -- SCRIBBLE / SKETCHBOOK EXPLAINER.
% Warm paper, faint dot-grid, wobbly ink outlines (random-steps deco),
% hand-curved arrows, highlighter swooshes, doodles. Comic Neue.
% Every slide teaches one idea, drawn by hand. Two pdflatex passes.
% =====================================================================

\usepackage[T1]{fontenc}
\usepackage[default]{comicneue}
\usepackage{textcomp}
\usepackage{pifont}
\usepackage{tikz}
\usetikzlibrary{decorations.pathmorphing, calc, arrows.meta, shapes.misc,
  positioning, backgrounds}

\usetheme{default}
\setbeamertemplate{navigation symbols}{}
\setbeamertemplate{footline}{}

\definecolor{paper}{HTML}{FAF4E6}
\definecolor{ink}{HTML}{34302A}
\definecolor{inkblue}{HTML}{2F5DA6}
\definecolor{mred}{HTML}{D14B4B}
\definecolor{myellow}{HTML}{F4CE5E}
\definecolor{mgreen}{HTML}{4CA362}
\definecolor{morange}{HTML}{E0883C}
\definecolor{mpurple}{HTML}{7A5BA6}
\definecolor{pencil}{HTML}{9A938A}

\setbeamercolor{background canvas}{bg=paper}
\setbeamercolor{normal text}{fg=ink}

\pgfmathsetseed{55}

\tikzset{
  hand/.style={decorate, decoration={random steps, segment length=5pt, amplitude=0.5pt}},
  pen/.style={hand, draw=ink, line width=1.2pt, line cap=round, line join=round},
}

% full-bleed sketch frame: (0,0)=SW .. (16,9)=NE
\newcommand{\art}[1]{%
  \begin{frame}[plain]
  \begin{tikzpicture}[remember picture, overlay,
      shift={(current page.south west)}, x=1cm, y=1cm]
    \fill[paper] (0,0) rectangle (16,9);
    \foreach \gx in {0.5,1.1,...,15.7}{\foreach \gy in {0.4,1.0,...,8.8}{
       \fill[ink, opacity=0.05] (\gx,\gy) circle (0.017);}}
    #1
  \end{tikzpicture}
  \end{frame}}

% hand-curved arrow with a clean tip
\newcommand{\harrow}[2]{%
  \draw[draw=ink, line width=1.1pt, line cap=round, -{Stealth[length=2.6mm]}]
     #1 to[bend left=5] #2;}
\newcommand{\harrowc}[3]{%
  \draw[draw=#3, line width=1.1pt, line cap=round, -{Stealth[length=2.6mm]}]
     #1 to[bend left=5] #2;}

% highlighter swoosh
\newcommand{\hl}[3]{\draw[#3, line width=7pt, opacity=0.35, line cap=round] #1 -- #2;}

% a sketchy box: {(c)}{halfw}{halfh}{fillcolor}
\newcommand{\dbox}[4]{%
  \draw[pen, fill=#4] ($#1+(-#2,-#3)$) rectangle ($#1+(#2,#3)$);}

% ---- doodles --------------------------------------------------------------
\newcommand{\magnifier}[2]{%
  \begin{scope}[shift={#1}, scale=#2]
    \draw[pen, fill=inkblue!8] (0,0) circle (0.55);
    \draw[pen, line width=2.2pt] (0.42,-0.42) -- (1.05,-1.05);
  \end{scope}}

\newcommand{\lockdoodle}[3]{% center open(0/1) color
  \begin{scope}[shift={#1}]
    \draw[pen, draw=#3, fill=#3!16] (-0.42,-0.48) rectangle (0.42,0.36);
    \ifnum#2=1 \draw[pen, draw=#3, line width=1.4pt] (-0.26,0.36) arc (205:35:0.30);
    \else      \draw[pen, draw=#3, line width=1.4pt] (-0.26,0.36) arc (180:0:0.26);\fi
    \fill[#3] (0,-0.02) circle (0.05);
    \draw[draw=#3, line width=1pt] (0,-0.02)--(0,-0.24);
  \end{scope}}

\newcommand{\crab}[2]{% Ferris the crab
  \begin{scope}[shift={#1}, scale=#2]
    \draw[pen, draw=morange!60!black, fill=morange!85] (0,0) ellipse (0.72 and 0.42);
    \draw[pen, draw=morange!60!black] (-0.18,0.34)--(-0.18,0.62);
    \draw[pen, draw=morange!60!black] (0.18,0.34)--(0.18,0.62);
    \fill[white] (-0.18,0.66) circle (0.1);\fill[white] (0.18,0.66) circle (0.1);
    \fill[ink] (-0.18,0.66) circle (0.04);\fill[ink] (0.18,0.66) circle (0.04);
    \draw[pen, draw=morange!70!black, line width=1.6pt] (-0.72,0.02) .. controls (-1.15,0.12) .. (-1.18,-0.26);
    \draw[pen, draw=morange!70!black, line width=1.6pt] (0.72,0.02) .. controls (1.15,0.12) .. (1.18,-0.26);
    \draw[pen, draw=morange!70!black] (-1.18,-0.26) .. controls (-1.0,-0.04) .. (-0.86,-0.2);
    \draw[pen, draw=morange!70!black] (1.18,-0.26) .. controls (1.0,-0.04) .. (0.86,-0.2);
    \foreach \i in {-0.34,-0.1,0.14}{
      \draw[pen, draw=morange!70!black, line width=1pt] (\i,-0.4)--(\i-0.12,-0.62);
      \draw[pen, draw=morange!70!black, line width=1pt] (\i+0.5,-0.4)--(\i+0.62,-0.62);}
    \draw[draw=ink, line width=1pt, line cap=round] (-0.12,0.06) .. controls (0,-0.05) .. (0.12,0.06);
  \end{scope}}

\newcommand{\rocket}[2]{% rocket pointing right
  \begin{scope}[shift={#1}, scale=#2]
    \draw[pen, draw=inkblue, fill=inkblue!12]
       (-0.6,-0.28) -- (0.5,-0.28) .. controls (0.95,-0.1) and (0.95,0.1) .. (0.5,0.28)
       -- (-0.6,0.28) -- cycle;
    \draw[pen, draw=inkblue] (0.18,0) circle (0.13);
    \draw[pen, draw=inkblue, fill=inkblue!22] (-0.6,0.28)--(-0.88,0.52)--(-0.45,0.28)--cycle;
    \draw[pen, draw=inkblue, fill=inkblue!22] (-0.6,-0.28)--(-0.88,-0.52)--(-0.45,-0.28)--cycle;
    \draw[pen, draw=morange, fill=myellow!60] (-0.6,-0.14)--(-1.05,0)--(-0.6,0.14)--cycle;
  \end{scope}}

\newcommand{\snail}[2]{%
  \begin{scope}[shift={#1}, scale=#2]
    \draw[pen, draw=mgreen!50!black, fill=mgreen!16]
       (-0.95,-0.32) .. controls (-1.15,0.1) .. (-0.72,0.1) -- (0.25,0.1)
       .. controls (0.45,0.1) .. (0.32,-0.32) -- cycle;
    \draw[pen, draw=mgreen!45!black, fill=mgreen!10] (0.02,0.02) circle (0.42);
    \draw[pen, draw=mgreen!45!black] (0.02,0.02) circle (0.26);
    \draw[pen, draw=mgreen!45!black] (0.02,0.02) circle (0.12);
    \draw[pen, draw=mgreen!50!black] (-0.95,-0.06)--(-1.1,0.24);
    \draw[pen, draw=mgreen!50!black] (-0.82,-0.06)--(-0.94,0.28);
    \fill[ink] (-1.1,0.26) circle (0.03);\fill[ink] (-0.94,0.30) circle (0.03);
  \end{scope}}

\newcommand{\bug}[3]{% center scale color
  \begin{scope}[shift={#1}, scale=#2]
    \draw[pen, draw=#3!55!black, fill=#3!80] (0,0) ellipse (0.27 and 0.37);
    \draw[pen, draw=#3!55!black] (0,0.37)--(0,-0.37);
    \foreach \y in {0.16,-0.02,-0.2}{
      \draw[pen, draw=#3!55!black, line width=1pt] (-0.27,\y)--(-0.5,\y+0.09);
      \draw[pen, draw=#3!55!black, line width=1pt] (0.27,\y)--(0.5,\y+0.09);}
    \draw[pen, draw=#3!55!black] (-0.1,0.37)--(-0.2,0.56);
    \draw[pen, draw=#3!55!black] (0.1,0.37)--(0.2,0.56);
  \end{scope}}

\newcommand{\tick}[2]{\draw[draw=mgreen, line width=2pt, line cap=round, line join=round]
   ($#1+(-0.16,0)$) -- ($#1+(-0.03,-0.15)$) -- ($#1+(0.22,0.2)$);}

\begin{document}

% =====================================================================
% 1. TITLE
% =====================================================================
\art{
  \node[ink, anchor=west] at (1.3,6.7) {\fontsize{45}{49}\selectfont\textbf{Rotor in Rust}};
  \hl{(1.5,6.15)}{(7.7,6.15)}{myellow}
  \node[ink, anchor=west, align=left, text width=9.2cm] at (1.4,5.0)
     {\large a hand-drawn tour of what we built\\ --- and why you can trust it};
  \node[pencil, anchor=west] at (1.4,3.7) {\normalsize Jasmin Begi\'c \;\textbar\; Daniel Wassie};
  % Ferris
  \crab{(13.4,5.5)}{1.45}
  \node[inkblue] at (13.4,3.9) {\small ``Ferris'' says hi};
  % a little gear/rotor doodle, bottom-left
  \draw[pen, draw=ink, fill=myellow!30] (2.4,2.2) circle (0.7);
  \foreach \a in {0,45,...,315}{\draw[pen, line width=1.4pt]
     ({2.4+0.7*cos(\a)},{2.2+0.7*sin(\a)}) -- ({2.4+0.95*cos(\a)},{2.2+0.95*sin(\a)});}
  \draw[pen, fill=paper] (2.4,2.2) circle (0.26);
  \harrow{(3.6,2.2)}{(4.7,2.2)}
  \node[ink, anchor=west] at (4.9,2.2) {\small a ``rotor'' turns code into something checkable};
}

% =====================================================================
% 2. THE PROBLEM
% =====================================================================
\art{
  \node[ink] at (8,8.3) {\fontsize{26}{28}\selectfont\textbf{you can't test your way to safe}};
  \hl{(4.5,7.85)}{(11.5,7.85)}{myellow}
  % a big bin of inputs
  \draw[pen, fill=inkblue!7] (1.6,1.4) -- (2.3,6.0) -- (7.7,6.0) -- (8.4,1.4) -- cycle;
  \node[inkblue] at (5.0,6.5) {\small all possible inputs};
  \foreach \i in {1,...,70}{
     \pgfmathsetmacro{\px}{2.4+rnd*5.2}\pgfmathsetmacro{\py}{1.7+rnd*3.9}
     \fill[pencil, opacity=0.6] (\px,\py) circle (0.04);}
  \node[ink] at (5.0,3.6) {\fontsize{34}{36}\selectfont\textbf{2\textsuperscript{64}}};
  % the hidden bug
  \bug{(6.9,2.3)}{0.8}{mred}
  \harrowc{(9.2,2.0)}{(7.4,2.2)}{mred}
  \node[mred, align=left] at (11.0,1.9) {\small one of them\\ crashes};
  % magnifier checking a few
  \magnifier{(10.6,5.2)}{1.2}
  \node[ink, align=left] at (12.7,5.6)
     {\small you try a\\ handful \dots\\ \textit{the bug hides}\\ \textit{in the rest}};
}

% =====================================================================
% 3. THE BIG IDEA
% =====================================================================
\art{
  \node[ink] at (8,8.3) {\fontsize{26}{28}\selectfont\textbf{so ask ONE question instead}};
  \hl{(4.6,7.85)}{(11.4,7.85)}{myellow}
  % thought bubble
  \draw[pen, fill=paper] (5.0,5.9) ellipse (3.6 and 1.5);
  \draw[pen, fill=paper] (2.1,4.3) circle (0.22);
  \draw[pen, fill=paper] (2.5,4.8) circle (0.32);
  \node[ink, align=center] at (5.0,5.9)
     {\large can it \emph{ever} reach a bad state,\\ for \textbf{any} input, within k steps?};
  % two outcomes
  \harrow{(8.7,5.9)}{(10.2,6.7)}
  \harrow{(8.7,5.6)}{(10.2,4.4)}
  \node[mgreen, align=left] at (12.8,6.8) {\textbf{YES} $\rightarrow$ here is the\\ exact input that does it};
  \node[inkblue, align=left] at (12.8,4.3) {\textbf{NO} $\rightarrow$ a real\\ guarantee, for all inputs};
  \node[ink] at (7.2,2.3) {\small treat the unknown input like \emph{algebra}, not like a guess};
  \hl{(3.0,2.25)}{(11.4,2.25)}{mgreen}
}

% =====================================================================
% 4. THE PIPELINE
% =====================================================================
\art{
  \node[ink] at (8,8.3) {\fontsize{26}{28}\selectfont\textbf{the assembly line}};
  \hl{(6.0,7.85)}{(10.0,7.85)}{myellow}
  \dbox{(2.3,5.2)}{1.2}{0.7}{pencil!12}
  \node[ink, align=center] at (2.3,5.2) {\small compiled\\ binary};
  \dbox{(6.4,5.2)}{1.2}{0.7}{myellow!35}
  \node[ink, align=center] at (6.4,5.2) {\small the\\ \textbf{model}};
  \dbox{(10.5,5.2)}{1.2}{0.7}{pencil!12}
  \node[ink, align=center] at (10.5,5.2) {\small the solver\\ (btormc)};
  \dbox{(14.0,5.2)}{1.1}{0.7}{myellow!35}
  \node[ink, align=center] at (14.0,5.2) {\small a picture\\ you read};
  \harrow{(3.6,5.2)}{(5.1,5.2)}
  \harrow{(7.7,5.2)}{(9.2,5.2)}
  \harrow{(11.8,5.2)}{(12.8,5.2)}
  \node[inkblue] at (4.35,5.9) {\scriptsize rotor};
  \node[inkblue] at (12.3,5.9) {\scriptsize visualiser};
  % "us"
  \draw[pen, draw=mred] (5.1,3.9) rectangle (7.7,3.0);
  \node[mred] at (6.4,3.45) {\small \textbf{we built these two}};
  \harrowc{(6.4,3.95)}{(6.4,4.45)}{mred}
  \harrowc{(13.0,4.2)}{(13.7,4.45)}{mred}
  \node[ink] at (8,1.6) {\small what's verified is the \emph{real machine code} --- exactly what runs};
}

% =====================================================================
% 5. WHAT THE MODEL IS
% =====================================================================
\art{
  \node[ink] at (8,8.3) {\fontsize{26}{28}\selectfont\textbf{what rotor actually draws}};
  \hl{(5.0,7.85)}{(11.0,7.85)}{myellow}
  % the machine
  \draw[pen, fill=inkblue!6] (1.4,2.2) rectangle (8.4,6.6);
  \node[inkblue] at (2.5,6.2) {\small the machine};
  \node[ink] at (2.5,5.5) {\scriptsize registers};
  \foreach \i in {0,...,5}{\draw[pen, fill=paper] (1.9+\i*0.7,4.9) rectangle (2.4+\i*0.7,5.3);}
  \node[ink] at (2.3,4.4) {\scriptsize memory};
  \foreach \i in {0,...,7}{\foreach \j in {0,1}{
     \draw[pen, fill=paper] (1.9+\i*0.55,3.2+\j*0.5) rectangle (2.3+\i*0.55,3.6+\j*0.5);}}
  \draw[pen, -{Stealth[length=2mm]}] (6.6,3.4) to[bend left=40] (6.6,4.6);
  \node[ink] at (7.4,4.0) {\scriptsize rules};
  % the checklist
  \draw[pen, fill=mred!6] (9.4,2.2) rectangle (15.0,6.6);
  \node[mred] at (12.2,6.15) {\small must \textbf{NEVER} happen};
  \foreach \y/\t in {5.3/{divide by zero}, 4.7/{read bad memory},
       4.1/{illegal instruction}, 3.5/{wrong exit code}}{
     \tick{(10.0,\y)}{1}\node[ink, anchor=west] at (10.4,\y) {\scriptsize \t};}
  \node[pencil, anchor=west] at (10.0,2.9) {\scriptsize \dots 24 of these in total};
  \node[ink] at (8,1.5) {\small a position, the legal moves, and the forbidden squares};
}

% =====================================================================
% 6. THE REWRITE
% =====================================================================
\art{
  \node[ink] at (8,8.3) {\fontsize{26}{28}\selectfont\textbf{step 1 --- tidy the kitchen}};
  \hl{(5.0,7.85)}{(11.0,7.85)}{myellow}
  % messy ball
  \begin{scope}
  \foreach \i in {1,...,60}{
     \pgfmathsetmacro{\a}{rnd*360}\pgfmathsetmacro{\r}{0.2+rnd*1.5}
     \pgfmathsetmacro{\b}{rnd*360}\pgfmathsetmacro{\s}{0.2+rnd*1.5}
     \draw[pencil, line width=0.6pt, opacity=0.7]
        ({3.6+\r*cos(\a)},{4.7+\r*sin(\a)}) .. controls (3.6,4.7) .. ({3.6+\s*cos(\b)},{4.7+\s*sin(\b)});}
  \end{scope}
  \node[ink, align=center] at (3.6,2.4) {\small \texttt{rotor.c}\\ 14{,}000 lines, \emph{one file}};
  \harrow{(5.9,4.7)}{(8.1,4.7)}
  \node[ink] at (7.0,5.3) {\small rebuild};
  % neat boxes
  \foreach \y/\n in {6.0/btor2, 5.0/riscv, 4.0/machine, 3.0/model}{
     \draw[pen, fill=mgreen!10] (9.6,\y-0.4) rectangle (13.4,\y+0.4);
     \node[ink] at (11.5,\y) {\small \n};}
  \node[mgreen, align=center] at (11.5,2.2) {\small four clean parts\\ (Rust)};
  \node[ink] at (8,1.3) {\small same job --- now you can actually \emph{read} it};
}

% =====================================================================
% 7. SPEED
% =====================================================================
\art{
  \node[ink] at (8,8.3) {\fontsize{26}{28}\selectfont\textbf{step 2 --- and it got FAST}};
  \hl{(5.4,7.85)}{(10.6,7.85)}{myellow}
  \snail{(3.6,5.2)}{1.8}
  \node[ink, align=center] at (3.6,3.2) {\small old C rotor\\ \textbf{139 seconds}};
  \rocket{(11.8,5.4)}{2.2}
  \node[ink, align=center] at (11.8,3.2) {\small Rust rotor\\ \textbf{0.1 seconds}};
  \harrow{(5.6,4.6)}{(9.4,4.9)}
  \node[ink] at (7.6,6.4) {\fontsize{40}{42}\selectfont\textbf{\textasciitilde 1000\texttimes}};
  \hl{(6.0,6.0)}{(9.2,6.0)}{morange}
  \node[ink] at (8,1.5) {\small 428 MB of memory $\rightarrow$ just 20 MB, too};
}

% =====================================================================
% 8. THE COUNT
% =====================================================================
\art{
  \node[ink] at (8,8.3) {\fontsize{26}{28}\selectfont\textbf{wait --- WHY? don't trust it, count}};
  \hl{(4.2,7.85)}{(11.8,7.85)}{myellow}
  \node[ink] at (8,7.0) {\small both tools keep asking: \emph{``have I already built this piece?''}};
  % left: long list walked
  \node[inkblue] at (3.8,6.2) {\small old C: walk the whole list};
  \draw[pen, fill=paper] (2.8,1.4) rectangle (4.8,5.6);
  \foreach \y in {1.7,2.1,...,5.3}{\draw[pencil, line width=0.6pt] (3.0,\y)--(4.6,\y);}
  \draw[pen, -{Stealth[length=2mm]}] (5.1,5.3) to[bend left=30] (5.1,1.8);
  \node[mred, align=center] at (3.8,0.9) {\small 11.7 BILLION\\ comparisons};
  % right: index lookup
  \node[inkblue] at (12.0,6.2) {\small Rust: one lookup};
  \draw[pen, fill=mgreen!8] (10.6,3.6) rectangle (13.4,5.4);
  \foreach \i in {0,1,2}{\draw[pen] (10.6,3.6+\i*0.6) rectangle (13.4,4.2+\i*0.6);}
  \fill[myellow!60] (11.0,4.2) rectangle (13.0,4.8);
  \draw[pen, -{Stealth[length=2.4mm]}] (12.0,6.0) -- (12.0,5.45);
  \node[mgreen, align=center] at (12.0,2.9) {\small 159 THOUSAND\\ probes};
  \node[ink] at (8,1.5) {\small same question, same answer --- a \emph{list} vs.\ an \emph{index}};
}

% =====================================================================
% 9. THREE ROTORS
% =====================================================================
\art{
  \node[ink] at (8,8.3) {\fontsize{26}{28}\selectfont\textbf{actually, there are THREE rotors}};
  \hl{(4.4,7.85)}{(11.6,7.85)}{myellow}
  \foreach \x/\fc/\name/\face in {%
     3.2/pencil!15/{Original C}/{},
     8.0/myellow!35/{C + the trick}/{},
     12.8/mgreen!18/{Rust}/{}}{
     \draw[pen, fill=\fc] (\x-1.7,3.6) rectangle (\x+1.7,6.4);
     \node[ink] at (\x,6.0) {\small \textbf{\name}};}
  % faces / marks
  \draw[pen] (3.2,5.2) circle (0.5);
  \draw[ink, line width=1.2pt] (3.0,5.0) .. controls (3.2,4.8) .. (3.4,5.0);
  \fill[ink] (3.02,5.3) circle (0.05);\fill[ink] (3.38,5.3) circle (0.05);
  \draw[pen, fill=myellow!50] (8,5.2) circle (0.5);
  \node at (8,5.2) {\small\textbf{!}};  % lightbulb idea
  \draw[pen, draw=morange] (8,5.75) -- (8,5.95);
  \draw[pen, fill=mgreen!40] (12.8,5.2) circle (0.5);
  \node at (12.8,5.2) {\small\textbf{\ding{72}}};
  % stats
  \node[ink, align=center] at (3.2,4.3) {\scriptsize 139 s $\cdot$ 428 MB\\ 11.7 B compares};
  \node[ink, align=center] at (8,4.3) {\scriptsize 1.5 s $\cdot$ 95\texttimes\\ \emph{byte-identical}};
  \node[ink, align=center] at (12.8,4.3) {\scriptsize 0.1 s $\cdot$ 20 MB\\ 1000\texttimes\ faster};
  \harrow{(5.0,5.0)}{(6.2,5.0)}
  \harrow{(9.8,5.0)}{(11.0,5.0)}
  \node[ink] at (8,2.4) {\small the speed was never the \emph{language} --- it was one small idea};
  \hl{(3.4,2.35)}{(12.6,2.35)}{mgreen}
  \node[pencil] at (8,1.4) {\small ``perfection is lots of little things done well''};
}

% =====================================================================
% 10. THE BACK-PORT
% =====================================================================
\art{
  \node[ink] at (8,8.3) {\fontsize{26}{28}\selectfont\textbf{we taught the old dog the trick}};
  \hl{(4.6,7.85)}{(11.4,7.85)}{myellow}
  \draw[pen, fill=mgreen!12] (2.0,4.0) rectangle (5.2,6.2);
  \node[ink, align=center] at (3.6,5.1) {\small Rust\\ \textbf{idea}};
  \draw[pen, fill=pencil!14] (10.8,4.0) rectangle (14.0,6.2);
  \node[ink, align=center] at (12.4,5.1) {\small the original\\ \textbf{C}};
  % carry a lightbulb/key across
  \harrow{(5.4,5.1)}{(10.6,5.1)}
  \node[ink] at (8,5.7) {\small drop in the hash};
  \draw[pen, fill=myellow!55] (8,4.4) circle (0.28);
  \node at (8,4.4) {\small\textbf{!}};
  \node[ink, align=center] at (8,2.7)
     {\large \textbf{95\texttimes\ faster} in C itself ---\\ and the output is \emph{byte-for-byte the same}};
  \hl{(4.4,1.7)}{(11.6,1.7)}{morange}
  \node[ink] at (8,1.5) {\small proof it's the idea, not the language};
}

% =====================================================================
% 11. SYMBOLIC ARGV
% =====================================================================
\art{
  \node[ink] at (8,8.3) {\fontsize{26}{28}\selectfont\textbf{opening a door that was nailed shut}};
  \hl{(4.0,7.85)}{(12.0,7.85)}{myellow}
  \node[ink] at (8,6.9) {\small programs also take input from the words after their name};
  \node[ink] at (3.0,5.6) {\Large \texttt{./prog}};
  % frozen args
  \lockdoodle{(5.4,5.5)}{0}{pencil}
  \lockdoodle{(6.7,5.5)}{0}{pencil}
  \node[pencil] at (6.05,4.6) {\small frozen};
  \harrow{(7.6,5.5)}{(9.0,5.5)}
  % open args with ?
  \lockdoodle{(9.8,5.5)}{1}{inkblue}\node[inkblue] at (9.8,5.5) {\scriptsize ?};
  \lockdoodle{(11.1,5.5)}{1}{mred}\node[mred] at (11.1,5.5) {\scriptsize ?};
  \lockdoodle{(12.4,5.5)}{1}{mgreen}\node[mgreen] at (12.4,5.5) {\scriptsize ?};
  \node[inkblue] at (11.1,4.6) {\small open};
  \node[ink, align=center] at (8,3.1)
     {\small now the solver can \emph{choose} those bytes --- and find bugs\\
      that were impossible to reach before};
  \node[pencil] at (8,1.6) {\small (this was the original paper's own ``future work'' --- we built it)};
}

% =====================================================================
% 12. THE X / Y SECRET
% =====================================================================
\art{
  \node[ink] at (8,8.3) {\fontsize{26}{28}\selectfont\textbf{cracking a secret password}};
  \hl{(5.0,7.85)}{(11.0,7.85)}{myellow}
  % a safe with a rule
  \draw[pen, fill=ink!4] (1.8,2.4) rectangle (7.0,6.4);
  \draw[pen] (5.8,4.4) circle (0.7);
  \foreach \a in {0,30,...,330}{\draw[pen, line width=1pt]
     ({5.8+0.7*cos(\a)},{4.4+0.7*sin(\a)}) -- ({5.8+0.85*cos(\a)},{4.4+0.85*sin(\a)});}
  \draw[pen, line width=1.6pt] (5.8,4.4)--(6.1,4.9);
  \node[ink, align=center] at (3.4,5.4) {\small crashes \emph{only} if\\ arg1 = \textbf{X}\\ \textbf{and} arg2 = \textbf{Y}};
  \harrow{(7.3,4.4)}{(9.0,4.4)}
  % the solver cracks it
  \node[mgreen, align=center] at (12.2,5.4)
     {\Large \textbf{got it.}};
  \node[ink, align=center] at (12.2,4.2)
     {\small arg1 starts \textbf{X}\\ arg2 starts \textbf{Y}};
  \node[ink] at (8,2.6) {\small two letters, in two places, that must meet at once};
  \hl{(4.4,2.55)}{(11.6,2.55)}{mgreen}
  \node[pencil] at (8,1.6) {\small no keyboard could ever trigger it --- the solver found it by reason};
}

% =====================================================================
% 13. DIFFERENTIAL VALIDATION
% =====================================================================
\art{
  \node[ink] at (8,8.3) {\fontsize{26}{28}\selectfont\textbf{but is the new one CORRECT?}};
  \hl{(4.6,7.85)}{(11.4,7.85)}{myellow}
  \node[ink] at (8,6.9) {\small let the solver be the referee};
  \draw[pen, fill=pencil!12] (1.6,4.2) rectangle (4.6,6.0);
  \node[ink, align=center] at (3.1,5.1) {\small old C\\ rotor};
  \draw[pen, fill=mgreen!12] (11.4,4.2) rectangle (14.4,6.0);
  \node[ink, align=center] at (12.9,5.1) {\small Rust\\ rotor};
  % both point to a referee
  \draw[pen, fill=myellow!30] (8,5.1) circle (0.95);
  \node[ink, align=center] at (8,5.1) {\small same\\ answer?};
  \harrow{(4.8,5.1)}{(7.0,5.1)}
  \harrowc{(11.2,5.1)}{(9.0,5.1)}{ink}
  \node[ink] at (8,3.5) {\small same bug, at the same step number --- a fingerprint of the whole run};
  % big tally
  \node[mgreen] at (8,2.4) {\fontsize{40}{42}\selectfont\textbf{36 / 36}};
  \tick{(5.6,2.4)}{1}\tick{(10.2,2.4)}{1}
  \node[pencil] at (8,1.4) {\small every benchmark, two settings, zero disagreements};
}

% =====================================================================
% 14. THE VISUALIZER
% =====================================================================
\art{
  \node[ink] at (8,8.3) {\fontsize{26}{28}\selectfont\textbf{make the answer readable}};
  \hl{(5.0,7.85)}{(11.0,7.85)}{myellow}
  % left: wall of binary + confusion
  \draw[pen, fill=paper] (1.4,2.4) rectangle (6.4,6.6);
  \foreach \y in {2.8,3.3,...,6.2}{\foreach \bx in {0,...,9}{
     \pgfmathtruncatemacro{\bit}{random(0,1)}
     \node[pencil, font=\ttfamily\scriptsize] at (1.9+\bx*0.42,\y) {\bit};}}
  \node[mred] at (3.9,7.0) {\small ?!?};
  \harrow{(6.8,4.5)}{(8.6,4.5)}
  % right: a friendly graph
  \coordinate (g0) at (12.4,6.2);
  \coordinate (g1) at (10.8,4.6);\coordinate (g2) at (14.0,4.6);
  \coordinate (g3) at (10.0,3.0);\coordinate (g4) at (12.0,3.0);
  \draw[pen] (g0)--(g1) (g0)--(g2) (g1)--(g3) (g1)--(g4);
  \draw[pen, fill=mred!25] (12.4,6.2) circle (0.42);
  \node[mred, font=\scriptsize\bfseries] at (12.4,6.2) {BAD};
  \foreach \g in {g1,g2,g3,g4}{\draw[pen, fill=inkblue!18] (\g) circle (0.34);}
  \node[ink] at (8,1.6) {\small the wall of binary becomes a picture you can \emph{point at}};
}

% =====================================================================
% 15. THE LESSON / THANKS
% =====================================================================
\art{
  \node[ink, align=center] at (8,6.4)
     {\fontsize{30}{36}\selectfont\textbf{the hard part wasn't making it fast.}};
  \node[ink, align=center] at (8,5.1)
     {\fontsize{30}{36}\selectfont\textbf{it was proving it still tells the truth.}};
  \hl{(3.2,4.75)}{(12.8,4.75)}{myellow}
  \crab{(8,3.0)}{1.7}
  \node[ink] at (8,1.7) {\Large \textbf{thank you!}};
  \node[pencil] at (8,1.0)
     {\small github.com/jaspek/rotor-rust \;\textbar\; rotor-c-hashcons \;\textbar\; jaspek.github.io/rotor-rust};
}

\end{document}
